% GENERATED FILE. Edit canonical lessons, then run npm run generate:books.
% canonical-source-hash: fnv1a64:0f457a5d00a8032a
% canonical-lessons: ZH-C90-hui, ZH-C90-neng, ZH-W90-neng, ZH-C90-keyi, ZH-C90-youyong, ZH-C90-practice

\chapter{What I Can Do}
\label{ch:zh-what-i-can-do}
\hlchaptermodality{90}

\begin{chapteropening}
\textbf{By the end of this chapter:} \emph{I can say what I can do and ask whether I may.}

Three ways to say can: \zh{会} for what you have learned, \zh{能} for what you are able or allowed to do, kěyǐ for asking permission, built on one new written character: \zh{能}.
\end{chapteropening}

\section[\texorpdfstring{\zh{会} (huì)}{会 (huì)}]{\zh{会} (huì) — can, having learned how}

\label{lesson:ZH-C90-hui}



\begin{quote}
A new chapter, and it opens with what you have learned to do.

\begin{itemize}
\raggedright
  \item \emph{Recall:} say \zh{给你} — \emph{gěi nǐ}, here you are
  \item \emph{Recall:} say \emph{qǐng zhàn zài zhèr}, please stand here — \textbf{R1}, the chapter you have finished
  \item \emph{Recall:} say \emph{sījī}, a driver — \textbf{R2}
\end{itemize}
\end{quote}

\subsection*{You'll want to know: \zh{会}}
\textbf{\zh{会}} — \emph{huì} — \textquotedblleft{}can (having learned how)\textquotedblright{}.

\begin{quote}
\textbf{\zh{我会开车}} \emph{wǒ huì kāichē} — I can drive.
\end{quote}

\begin{quote}
\textbf{\zh{我会说汉语}} \emph{wǒ huì shuō Hànyǔ} — I can speak Chinese.
\end{quote}

\begin{cousinweb}
\textbf{\zh{会}} (\emph{huì}) is inherited from Middle Chinese \emph{hwajH} (Guangyun, in Baxter's transcription), glossed 'to come together'. The sense 'can' is not given there, so how it got it is not settled.
\end{cousinweb}

\begin{grammarlens}[title={What you've built}]
\zh{会} + a verb says you have learned to do it: driving, a language, a skill. A \emph{sījī huì kāichē}: a driver can drive. You have written \zh{会} since \zh{一会儿}; this is another word in the same character.
\end{grammarlens}

\subsection*{Your turn}
\begin{itemize}
\raggedright
  \item \emph{Say it:} \emph{huì}
  \item \emph{Say it:} \emph{huì}, once more
  \item \emph{Say it:} \zh{我会开车} — \emph{wǒ huì kāichē} — I can drive.
  \item \emph{Say it:} \zh{我会说汉语} — \emph{wǒ huì shuō Hànyǔ} — I can speak Chinese.
  \item \emph{Recall:} say \emph{sījī}, then say a driver can drive: \emph{sījī huì kāichē}
\end{itemize}

\begin{tcolorbox}[breakable,colback=teal!4,colframe=teal!35!black,title={Before you move on}]
What does \zh{会} say? (\textbf{That you have learned how}.) Say you can drive. (\textbf{\zh{我会开车}}.)
\end{tcolorbox}
\section[\texorpdfstring{\zh{能} (néng)}{能 (néng)}]{\zh{能} (néng) — can; able, allowed}

\label{lesson:ZH-C90-neng}



\begin{quote}
Before the new one: last lesson's word, then two words about being ill.

\begin{itemize}
\raggedright
  \item \emph{Recall:} say \zh{我会开车} — \emph{wǒ huì kāichē}
  \item \emph{Recall:} say \emph{bìng}, to fall ill — \textbf{R4}, recognition at distance
  \item \emph{Recall:} say \emph{yīshēng}, a doctor — \textbf{R2}
\end{itemize}
\end{quote}

\subsection*{You'll want to know: \zh{能}}
\textbf{\zh{能}} — \emph{néng} — \textquotedblleft{}can; to be able to; to be allowed to\textquotedblright{}.

\begin{quote}
\textbf{\zh{我能去吗}} \emph{wǒ néng qù ma} — Can I go?
\end{quote}

\begin{quote}
\emph{wǒ bìng le, bù néng qù} — I'm ill, I can't go.
\end{quote}

The next lesson puts the new character in your hand.

\begin{cousinweb}
\textbf{\zh{能}} (\emph{néng}) is inherited from Middle Chinese \emph{nong} (Guangyun, in Baxter's transcription), glossed 'skilled; good at work'.
\end{cousinweb}

\begin{grammarlens}[title={What you've built}]
\zh{会} is a skill you have learned. \zh{能} is being able right now, or being allowed: you \zh{会} drive, but ill you \zh{不能}.
\end{grammarlens}

\subsection*{Your turn}
\begin{itemize}
\raggedright
  \item \emph{Say it:} \emph{néng}
  \item \emph{Say it:} \emph{néng}, once more
  \item \emph{Say it:} \zh{我能去吗} — \emph{wǒ néng qù ma} — Can I go?
  \item \emph{Say it:} \emph{wǒ bìng le, bù néng qù} — I'm ill, I can't go.
  \item \emph{Recall:} say \emph{yīshēng}, then \emph{néng} again
\end{itemize}

\begin{tcolorbox}[breakable,colback=teal!4,colframe=teal!35!black,title={Before you move on}]
What does \zh{能} say? (\textbf{Able now, or allowed}.) And \zh{会}? (\textbf{Learned how}.)
\end{tcolorbox}
\section[\texorpdfstring{\zh{能} (néng)}{能 (néng)}]{\zh{能} — \zh{厶} over the moon piece, then two \zh{匕}}

\label{lesson:ZH-W90-neng}



\begin{quote}
The word first, then the character inside it.

\begin{itemize}
\raggedright
  \item \emph{Recall:} say \emph{néng} — \textbf{R1}, one lesson back, before it decays
\end{itemize}

\begin{quote}
\zh{明} \zh{前}
\end{quote}

\emph{Read it:} each one aloud, with nothing said first — \emph{míng}, \emph{qián}
\end{quote}

\subsection*{Script — \zh{厶} over the moon piece, then two \zh{匕}}
\hlblockfigure{\detokenize{figures/ZH-W90-neng-filmstrip.pdf}}{How \zh{能} is written, stroke by stroke}

\begin{quote}
\zh{能}
\end{quote}

\textbf{The whole left side first, top to bottom, then the right.}

\textbf{\zh{厶}, two strokes}: a stroke that falls to the left and turns right, then a dot.

\textbf{The moon piece \zh{月}, four strokes}, as it stands in \zh{明} and \zh{前}: the left side, the top that turns down and hooks, two short bars inside.

\textbf{Two \zh{匕}, two strokes each}, the upper one first: a short left-falling stroke, then a vertical that bends right along its foot and hooks up.

2 + 4 + 4 = ten.

\subsection*{Your turn}
\begin{itemize}
\raggedright
  \item \emph{Write it:} \zh{能}, all ten, with the model in front of you
  \item \emph{Cover:} write it again from memory, and count the pieces, not the strokes
\end{itemize}

\begin{tcolorbox}[breakable,colback=teal!4,colframe=teal!35!black,title={Before you move on}]
How many strokes in \zh{能}? (\textbf{Ten}.) Which piece do you know from \zh{明}? (\textbf{The moon piece, \zh{月}}.) Which side is written first? (\textbf{The left, all of it}.)

Source: \href{https://raw.githubusercontent.com/chanind/hanzi-writer-data/68d10a4b21150cae5e1ebbd223eed289cf32d90c/data/\%E8\%83\%BD.json}{Hanzi Writer Data \zh{能}, pinned stroke paths and medians}.
\end{tcolorbox}
\section[kěyǐ]{kěyǐ — may; can, with permission}

\label{lesson:ZH-C90-keyi}



\begin{quote}
Before the new one: read the last lesson's character, then two words you will need with this one.

\begin{itemize}
\raggedright
  \item \emph{Read it:} \zh{我能去吗}, aloud — \emph{wǒ néng qù ma}
  \item \emph{Recall:} say \emph{zuò}, to sit — \textbf{R4}, recognition at distance
  \item \emph{Recall:} say \emph{gōngrén}, a worker — \textbf{R2}
\end{itemize}
\end{quote}

\subsection*{You'll want to know: kěyǐ}
\emph{kěyǐ} — \textquotedblleft{}may; can (with permission)\textquotedblright{}.

\begin{quote}
\emph{wǒ kěyǐ zuò zhèr ma} — May I sit here?
\end{quote}

Two dips in a row: spoken, the first one rises, \emph{kéyǐ}. It is written \emph{kěyǐ}.

Its characters come later; for now it is a word for the ear.

\begin{cousinweb}
\emph{kěyǐ} is built from kě 'may' and yǐ 'to use', both inherited from Middle Chinese, \emph{khaX} and \emph{yiX} (Guangyun, in Baxter's transcription), glossed 'to permit' and 'to use'.
\end{cousinweb}

\begin{grammarlens}[title={What you've built}]
\zh{会} learned, \zh{能} able or allowed, \emph{kěyǐ} permitted. Asking whether you may, \emph{kěyǐ} is the polite one; the answer is \emph{kěyǐ}, or \emph{bù kěyǐ}.
\end{grammarlens}

\subsection*{Your turn}
\begin{itemize}
\raggedright
  \item \emph{Say it:} \emph{kěyǐ} — and hear it as \emph{kéyǐ}
  \item \emph{Say it:} \emph{kěyǐ}, once more
  \item \emph{Say it:} \emph{wǒ kěyǐ zuò zhèr ma} — May I sit here?
  \item \emph{Say it:} \emph{kěyǐ} — yes, you may
  \item \emph{Recall:} say \emph{gōngrén}, then \emph{kěyǐ} again
\end{itemize}

\begin{tcolorbox}[breakable,colback=teal!4,colframe=teal!35!black,title={Before you move on}]
What does \emph{kěyǐ} ask? (\textbf{Whether you may}.) How is it said? (\emph{\textbf{kéyǐ}}, the first dip rising.)
\end{tcolorbox}
\section[yóuyǒng]{yóuyǒng — to swim}

\label{lesson:ZH-C90-youyong}



\begin{quote}
Before the new one: last lesson's word, then two kinds of exercise.

\begin{itemize}
\raggedright
  \item \emph{Recall:} say \emph{kěyǐ}, may — \textbf{R1}, one lesson back
  \item \emph{Recall:} say \emph{yùndòng}, sport — \textbf{R4}, recognition at distance
  \item \emph{Recall:} say \emph{pǎobù}, to run — \textbf{R4}
\end{itemize}
\end{quote}

\subsection*{You'll want to know: yóuyǒng}
\emph{yóuyǒng} — \textquotedblleft{}to swim\textquotedblright{}.

\begin{quote}
\emph{nǐ huì yóuyǒng ma} — Can you swim?
\end{quote}

Its characters come later; for now it is a word for the ear.

\begin{cousinweb}
\emph{yóuyǒng} is built from yóu 'to float' and yǒng 'to swim under water', both inherited from Middle Chinese, \emph{yuw} and \emph{hjwaengH} (Guangyun, in Baxter's transcription). \emph{hjwaengH} would regularly give yòng, with a fall; why the modern word dips is not settled.
\end{cousinweb}

\begin{grammarlens}[title={What you've built}]
Swimming is learned, so it takes \zh{会}: \emph{nǐ huì yóuyǒng ma}. \zh{能} asks whether you are able today: \emph{jīntiān nǐ néng yóuyǒng ma}.
\end{grammarlens}

\subsection*{Your turn}
\begin{itemize}
\raggedright
  \item \emph{Say it:} \emph{yóuyǒng}
  \item \emph{Say it:} \emph{yóuyǒng}, once more
  \item \emph{Say it:} \emph{nǐ huì yóuyǒng ma} — Can you swim?
  \item \emph{Say it:} \emph{wǒ bú huì yóuyǒng} — I can't swim.
  \item \emph{Recall:} say \emph{pǎobù}, then \emph{yóuyǒng} again
\end{itemize}

\begin{tcolorbox}[breakable,colback=teal!4,colframe=teal!35!black,title={Before you move on}]
What does \emph{yóuyǒng} mean? (\textquotedblleft{}To swim\textquotedblright{}.) Which can goes with it, \zh{会} or \zh{能}? (\textbf{\zh{会}, a skill you learn}.)
\end{tcolorbox}
\section[Practice]{\zh{我会开车}}

\label{lesson:ZH-C90-practice}



\begin{quote}
The chapter closes here, with three ways to say can in one conversation.

\begin{itemize}
\raggedright
  \item \emph{Recall:} say \emph{yóuyǒng}, to swim — \textbf{R1}, one lesson back
  \item \emph{Recall:} say \emph{shàngbān}, to go to work — \textbf{R2}
  \item \emph{Recall:} say \emph{gōngsī}, a company — \textbf{R2}
\end{itemize}
\end{quote}

\subsection*{The exchange}
\begin{quote}
\textbf{A:} \zh{你会开车吗}
\end{quote}

\begin{quote}
\emph{nǐ huì kāichē ma} — Can you drive?
\end{quote}

\begin{quote}
\textbf{B:} \zh{会}
\end{quote}

\begin{quote}
\emph{huì} — I can.
\end{quote}

\begin{quote}
\textbf{A:} \zh{你明天能开车吗}
\end{quote}

\begin{quote}
\emph{nǐ míngtiān néng kāichē ma} — Can you drive tomorrow?
\end{quote}

\begin{quote}
\textbf{B:} \emph{bù néng, wǒ míngtiān shàngbān} — No, I'm at work tomorrow.
\end{quote}

\begin{quote}
\textbf{A:} \emph{nǐ huì yóuyǒng ma} — Can you swim?
\end{quote}

\begin{quote}
\textbf{B:} \zh{不会}
\end{quote}

\begin{quote}
\emph{bú huì} — I can't.
\end{quote}

\begin{quote}
\textbf{A:} \emph{wǒ kěyǐ zuò zhèr ma} — May I sit here?
\end{quote}

\begin{quote}
\textbf{B:} \emph{kěyǐ} — Yes, please.
\end{quote}

The answer to each is the verb alone: \zh{会}, \zh{不能}, \zh{不会}, \emph{kěyǐ}. Mandarin answers yes or no by saying the verb again.

\begin{tcolorbox}[breakable,skin=enhanced,colback=orange!6,colframe=orange!45!black,boxrule=0.5pt,arc=1mm,left=6pt,right=6pt,top=4pt,bottom=4pt,fonttitle=\bfseries,title={Read this}]
\emph{Read it:} the exchange through once without stopping, then answer the question below
\end{tcolorbox}

\begin{grammarlens}[title={What you've built — three cans}]
\begin{itemize}
\raggedright
  \item \textbf{\zh{我会开车}} \emph{wǒ huì kāichē} — I have learned to drive.
  \item \textbf{\zh{我明天不能开车}} \emph{wǒ míngtiān bù néng kāichē} — tomorrow I'm not able to.
  \item \emph{wǒ kěyǐ zuò zhèr ma} — may I sit here?
\end{itemize}

\textbf{Listen.} \emph{Nǐ huì shuō Hànyǔ ma?}

\textbf{Speak.} Answer it, then say one thing you cannot do today, and ask permission for something.

\emph{Write it:} \zh{能} — the whole left side first, then the two hooked pieces on the right
\end{grammarlens}

\subsection*{Your turn}
\begin{itemize}
\raggedright
  \item \emph{Say it:} \zh{你会开车吗} — \emph{nǐ huì kāichē ma} — Can you drive?
  \item \emph{Say it:} \zh{你明天能开车吗} — \emph{nǐ míngtiān néng kāichē ma} — Can you drive tomorrow?
  \item \emph{Say it:} \emph{wǒ kěyǐ zuò zhèr ma} — May I sit here?
  \item \emph{Say it:} the whole exchange, both parts
  \item \emph{Write it:} \zh{我能去吗}
\end{itemize}

\begin{tcolorbox}[breakable,colback=teal!4,colframe=teal!35!black,title={Before you move on}]
Say you can drive. (\textbf{\zh{我会开车}}.) Ask whether you can go. (\textbf{\zh{我能去吗}}.) Ask whether you may sit here. (\emph{wǒ kěyǐ zuò zhèr ma}.) Which can is a skill? (\textbf{\zh{会}}.)
\end{tcolorbox}
